%% ============================================================
%%  chapters/chapter2.tex  –  Literature Review and Related Work
%% ============================================================
\chapter{Literature Review and Related Work}
\label{chap:literature}

This chapter synthesizes relevant research across four intersecting domains: behavioral biometrics and keystroke dynamics, learning analytics and educational trace analysis, online assessment and academic integrity, and deep learning architectures for sequence modeling. The goal is to position the proposed integrated module within the broader research landscape and articulate the specific research gap it addresses.

%% ── 2.1 Keystroke Dynamics and Behavioral Biometrics ──────────
\section{Keystroke Dynamics and Behavioral Biometrics}
\label{sec:keystroke-dynamics}

Behavioral biometrics leverages characteristic patterns of human behavior to authenticate identity. Unlike anatomical biometrics (fingerprint, iris), behavioral biometrics do not require special hardware and can be collected unobtrusively during regular interaction with digital systems.

\subsection{Foundational keystroke dynamics research}

The field of keystroke dynamics for authentication was pioneered by \citet{monrose2000keystroke}, who established that keystroke timing patterns can serve as a unique identity signature. Their seminal work introduced the quantification of dwell time (the time a key is held down), flight time (the time between successive keystrokes), and inter-key latencies as features for authentication. They showed that a set of typing-based features, when aggregated and compared using distance metrics, could achieve useful authentication separation between users with relatively low false rejection and false acceptance rates.

Expanding on this foundation, \citet{dhakal2018observations} conducted a massive-scale empirical study analyzing 136 million keystrokes collected through an instrumented text-entry system. Their findings revealed significant individual variability in typing patterns, including systematic differences in dwell time distributions, burst-typing behavior, and error-correction patterns. Importantly, their work documented that typing behavior is both stable (within-user consistency) and variable across contexts (between-user variance), motivating the need for context-adaptive authentication.

\subsection{Keystroke dynamics in online examination contexts}

Within educational contexts, keystroke authentication has been specifically investigated for online examination integrity. \citet{mattsson2020keystroke} studied keystroke dynamics as a continuous authentication mechanism during remote examinations. Their work demonstrated that keystroke-based identity scores could flag deviations from enrolled baselines when the typing pattern diverges significantly, providing a mechanism for suspicious session detection.

More recently, \citet{senarath2023reevaluating} conducted a comprehensive re-evaluation of keystroke dynamics for continuous authentication in a modern online assessment setting. They implemented a system that continuously monitors typing patterns during examinations and compared various distance metrics and similarity thresholds. Their findings indicated that continuous keystroke monitoring significantly improves the ability to detect imposters mid-session, compared to one-time login authentication. They also identified practical challenges: device switching, keyboard changes, and stress-induced behavioral variance can degrade performance if not accounted for.

\subsection{Deep learning approaches to keystroke biometrics}

Traditional keystroke authentication relied on hand-crafted statistical features (means, standard deviations, percentiles of timing distributions). While effective to a degree, these approaches struggled with context variability. The emergence of deep learning has transformed the field by learning automatic feature representations from raw keystroke sequences.

\citet{acien2022typenet} introduced TypeNet, a deep LSTM-based architecture specifically designed for keystroke biometrics. TypeNet processes sequences of keystroke timing vectors (typically 70 keystrokes per sequence) through two stacked LSTM layers with batch normalization and dropout for regularization. The output is a 128-dimensional embedding that captures the essence of a keystroke sequence. TypeNet embeddings have been shown to generalize well across different text domains and to produce more stable cross-context matching than traditional features.

The TypeNet architecture is motivated by successes in face recognition and speaker identification, where embedding-based approaches (e.g., FaceNet, VoxCeleb) have achieved state-of-the-art results. By analogy, TypeNet treats each keystroke sequence as a "sample" in an embedding space, and users are defined by their cluster of embeddings. Authentication is performed by comparing a new sequence embedding against the user's cached mean and standard deviation.

%% ── 2.2 Learning Analytics and Educational Trace Analysis ────
\section{Learning Analytics and Educational Trace Analysis}
\label{sec:learning-analytics}

Learning analytics is the field of mining and interpreting educational data to support student success and institutional improvement. A key sub-domain is the analysis of learner interaction traces (logs, event sequences) to infer cognitive and metacognitive processes.

\subsection{Keystroke and interaction traces in educational settings}

Programming education has been a natural laboratory for trace analysis because code editors, IDEs, and online judges generate rich, structured interaction logs. Researchers have demonstrated that keystroke-level and code-level traces reveal significant pedagogical information.

\citet{arakawa2021situ} investigated the detection of self-regulated learning (SRL) struggles in programming assignments. They mined keystroke and code-edit traces to identify "struggle points"—moments where students encounter difficulty and engage in troubleshooting. Their analysis revealed that struggle is characterized by increased deletion rates, rapid back-and-forth edits, and non-linear code-building patterns. Notably, struggle *identification* enables instructors to intervene earlier and provide targeted support, improving learning outcomes.

\citet{dong2021trace} extended this line of work by proposing methods to automatically identify meaningful progress and regression in code solutions from trace logs. They showed that quantitative metrics derived from keystroke and code-edit traces (such as code length over time, compilation success rate, and test coverage evolution) correlate with learning outcomes. This supports the claim that process-level indicators are pedagogically relevant proxies for learning.

\subsection{Cognitive load and mental effort estimation}

Keystroke and interaction patterns can also serve as proxies for cognitive load and mental effort. Pauses between keystrokes, for instance, often reflect thinking time or problem decomposition. High keystroke bursts may reflect confidence or automated code generation (or copy-pasting).

While not focusing narrowly on keystroke analysis, the broader literature on cognitive load theory \citep{sweller1988cognitive} suggests that learners exhibit measurable behavioral changes under different levels of cognitive load. Eye-tracking, typing speed, and pause patterns have all been used as indirect measures of cognitive load in experimental studies.

\subsection{Formative feedback and LLM-assisted intervention}

\citet{nguyen2023gpt4} demonstrated that large language models (LLMs) can generate tiered, personalized formative feedback on student code submissions. Their system uses code content, error messages, and submission context to prompt GPT-4, which generates feedback at multiple levels of abstraction (syntax, logic, style, conceptual). While their focus was on feedback content, not behavioral analysis, the work highlights the growing integration of LLMs in educational support systems.

In the proposed module, LLM integration is envisioned for deep behavioral analysis: given keystroke traces, deletion patterns, and code snapshots, an LLM can generate naturalistic, contextual explanations of inferred cognitive processes and suggest pedagogically appropriate interventions.

%% ── 2.3 Online Assessment and Academic Integrity ───────────────
\section{Online Assessment and Academic Integrity}
\label{sec:online-assessment}

The surge in online and blended education has escalated concerns about academic integrity in remote assessments. Traditional in-person invigilation is impractical or impossible. Consequently, institutions have deployed diverse technical and procedural controls.

\subsection{Remote proctoring and biometric verification}

Remote proctoring typically combines webcam surveillance, browser lockdown, and identity verification. \citet{rapchak2020remote} surveyed remote proctoring technologies and reviewed their effectiveness and privacy implications. While webcam-based approaches are common, they are expensive, intrusive, and raise equity concerns (e.g., privacy in shared living spaces, camera requirements for students with disabilities).

Keystroke dynamics offers a non-invasive alternative or complement to webcam proctoring: once enrolled, a student's typing pattern can be continuously monitored without additional hardware, privacy concerns, or accessibility barriers.

\subsection{Plagiarism and external assistance detection}

Static plagiarism detection tools (e.g., MOSS, JPlag, Turnitin) compare submitted code against a database of known solutions. However, they miss plagiarism from external AI services or custom one-off solutions. \citet{chaski2012forensic} applied forensic linguistics to detect stylistic inconsistencies in text, a technique that could be extended to code.

Keystroke dynamics can complement these approaches: anomalous typing patterns (excessive paste, unrealistic burst speeds, sudden stylistic shifts) can flag sessions meriting closer manual inspection.

\subsection{Ethical and regulatory considerations}

The deployment of biometric and behavioral analytics in educational settings raises privacy, fairness, and consent concerns. Institutions must:

\begin{enumerate}
  \item \textbf{Obtain informed consent:} Students must understand what is collected, how it is used, and for how long it is retained.
  \item \textbf{Ensure due process:} High-risk flags should trigger human review and student opportunity to respond.
  \item \textbf{Calibrate thresholds fairly:} Authentication and behavior thresholds should be calibrated to avoid disparate impact on students from different backgrounds (e.g., non-native English speakers, neurodivergent learners, users with motor disabilities).
  \item \textbf{Minimize data retention:} Keystroke data should not be retained indefinitely; institutions should define retention policies aligned with FERPA, GDPR, or equivalent regulations.
\end{enumerate}

%% ── 2.4 Deep Learning for Sequence Modeling ──────────────────
\section{Deep Learning for Sequence Modeling}
\label{sec:deep-learning}

The keystroke authentication and behavior analytics proposed in this dissertation rely on deep learning for sequence modeling. A brief review of foundational architectures is warranted.

\subsection{Recurrent Neural Networks and LSTM}

Recurrent Neural Networks (RNNs) were designed to process sequences by maintaining a hidden state that is updated at each timestep. However, vanilla RNNs suffer from vanishing and exploding gradients, limiting their ability to capture long-range dependencies.

\citet{hochreiter1997lstm} introduced Long Short-Term Memory (LSTM) networks, which augment RNNs with gating mechanisms that allow the network to selectively forget and remember information. LSTMs have become a standard architecture for sequence processing, including time-series analysis, natural language processing, and anomaly detection.

In the context of keystroke sequences, LSTM networks are well-suited because:
\begin{itemize}
  \item Keystroke sequences have variable length and temporal structure.
  \item Typing patterns are influenced by recent history (e.g., error correction ripples) but also by longer-term style.
  \item LSTM cells naturally capture both short-term bursts and long-term behavioral trends.
\end{itemize}

\subsection{Embedding-based authentication}

Instead of training LSTMs to output a classification (user A vs. user B vs. ...),modern biometric systems (TypeNet, FaceNet, etc.) train neural networks to output fixed-size embeddings that cluster samples from the same user and separate samples from different users. At inference time, a new sample's embedding is compared to reference embeddings using distance or similarity metrics (Euclidean distance, cosine similarity, etc.).

This embedding-based approach offers several advantages:
\begin{enumerate}
  \item \textbf{Scalability:} New users can be enrolled without retraining; only their reference embeddings are stored.
  \item \textbf{Interpretability:} Shared embedding space enables visualization and manual similarity inspection.
  \item \textbf{Flexibility:} Multiple users can be compared, enabling identification (1:N matching) in addition to verification (1:1 matching).
\end{enumerate}

\subsection{Convolutional Neural Networks and temporal convolutions}

While LSTMs are powerful, Convolutional Neural Networks (CNNs) have also been explored for sequence processing. CNNs apply local convolutional filters to capture patterns of fixed size, and hierarchical pooling reduces the sequence length. For keystroke sequences, CNNs might capture local digraph or trigraph timing patterns.

However, LSTMs are the more standard choice for keystroke dynamics in the current literature, partly because they more naturally capture variable-length sequences and long-range dependencies.

%% ── 2.5 Related Systems and Comparative Analysis ─────────────
\section{Related Systems and Comparative Analysis}
\label{sec:related-systems}

Several systems and research projects have combined authentication and learning analytics or addressed related multi-objective design problems. A brief review:

\subsection{Plagiarism detection systems with behavioral signals}

Systems like Turnitin and MOSS primarily analyze code similarity but have begun integrating behavioral signals (submission patterns, revision history). However, keystroke-level analysis is not standard in these tools.

\subsection{IDE-integrated analytics}

Tools like CodeTracer and Visual Studio's CodeLens collect IDE-level events (compilation, debugging, testing) alongside code edits. These provide richer semantic context than keystroke events alone but require IDE integration and are specific to certain programming environments.

\subsection{Educational multimodal analytics}

Some research platforms (e.g., with eye-tracking + keystroke + code traces) study multimodal signals. However, practical educational deployments rarely integrate multiple modalities due to cost and complexity. Keystroke alone is attractive because it is universal across devices and contexts.

\subsection{Positioning of this work}

The proposed module differs from existing work in several ways:

\begin{enumerate}
  \item \textbf{Integration:} Most systems separate authentication from analytics. This module aims to integrate both within one service boundary and event stream.
  
  \item \textbf{Interpretability focus:} Many biometric systems output scores or probabilities. This module emphasizes interpretable behavioral indicators (friction points, cognitive load estimates, engagement profiles) designed to inform instructor decision-making.
  
  \item \textbf{Deployment orientation:} The module is conceived as a production microservice with explicit APIs, database integration, and LMS integration patterns, not a research prototype.
  
  \item \textbf{Educational specificity:} Design choices (feature sets, scoring rubrics, output formats) are informed by educational workflows and stakeholder feedback, not pure algorithmic optimization.
\end{enumerate}

%% ── 2.6 Gaps and Research Opportunities ────────────────────────
\section{Gaps and Research Opportunities}
\label{sec:research-gaps}

Despite the maturity of keystroke dynamics and learning analytics as separate fields, several gaps remain:

\begin{enumerate}
  \item \textbf{Integration gap:} Few systems simultaneously and transparently integrate keystroke authentication with pedagogically meaningful behavior analytics.
  
  \item \textbf{Context adaptation gap:} Most keystroke systems are context-aware in design but lack practical mechanisms for dynamic threshold calibration as students type on different devices or under stress.
  
  \item \textbf{Fairness gap:} Limited public research on bias and fairness in keystroke biometrics across learners with diverse language backgrounds, motor abilities, keyboard experience, and neurodiversity.
  
  \item \textbf{Deployment gap:} Research prototypes rarely provide production-ready microservice architectures suitable for adoption in real LMS deployments.
  
  \item \textbf{Multimodal fusion gap:} Most systems focus on keystroke alone. Fusion with mouse telemetry, IDE events, and voice (if speaking-alouds are recorded) could enrich analysis but requires careful coordination.
\end{enumerate}

This dissertation aims to partially address the integration, deployment, and interpretability gaps, laying groundwork for future work on context adaptation and fairness.

\clearpage

